Úvodní poznámka. Toto oddílové shrnutí se zaměřuje na praktická doporučení a zásady pro udržení akademické integrity v době dostupných generativních nástrojů. Hlavní doporučení níže jsou uvedena jako praktické návrhy pro garanty předmětů a učitele; tam, kde jde o širší právní či institucionální závazky (DPIA, DPA, smluvní podmínky), odkazujte na příslušné interní postupy a právní konzultace (viz kapitoly 2.2 a 8). Některá doporučení jsou shrnuta jako obecné (general background knowledge).

Cíle oddílu
\begin{itemize}
  \item Pomoci garantům kategorizovat úlohy podle pravidelného nasazení AI a očekávané úrovně deklarace ze strany studentů (praktická politika).
  \item Návrh konkrétních postupů pro navržení úloh a workflow hodnocení, které snižují riziko nevhodného použití AI.
  \item Příklady formulací do sylabů a šablony pro pole deklarace použití AI při odevzdání.
\end{itemize}

Doporučená kategorizace úloh (pro zařazení do sylabu) — návrh (general background knowledge)
\begin{enumerate}
  \item Zakázáno: úlohy, u kterých použití AI by zásadním způsobem podkopalo hodnotu výstupu (např. krátké dozorované písemky, časově omezené in‑class writings). V sylabu uveďte jasné sankce a proces řešení podezření.
  \item Povoleno s deklarací: úlohy, kde je možné AI použít, ale vyžadujeme explicitní deklaraci, popis kroků a připojení procesních důkazů (drafty, verze, poznámky k práci) (general background knowledge).
  \item Vyžadováno využití AI (AI‑enhanced assignments): úlohy, které explicitně integrují AI jako nástroj — student musí popsat, jak AI použil, jaký prompt použil a doložit vlastní přidanou hodnotu (analýza, kritika, adaptace).
\end{enumerate}

Praktické konstrukce zadání a ověřitelnost (general background knowledge)
\begin{itemize}
  \item Preferujte procesně orientované body hodnocení (např. milníky, průběžné odevzdání částí, komentované verze), nikoli pouze konečný produkt.
  \item Požadujte artefakty, které je těžší získat z externího generátoru bez vlastní práce: pracovní poznámky, postupy řešení, lokální logy z vývojového repozitáře, screencasty vývoje nebo sběr původních dat.
  \item U technických úloh (např. matematika, programování) zařaďte požadavky na ukázku postupu — v některých nástrojích lze vyžadovat nahrání postupu řešení jako součást hodnocení (viz příklad z praxe v rámci školních nástrojů: povinnost nahrání postupu v aplikaci „Pokrok v matematice“) \cite{kb_ai_101_tipu_pdf_04e4a115_page_8_1}.
  \item Využívejte živé či polosimulované ústní ověření (spot‑checky, krátké obhajoby), zvláště u závěrečných prací či projektů, kde autenticita výstupu je klíčová (general background knowledge).
\end{itemize}

Formulace do sylabu — ukázky (copy‑\\ \&‑paste připravené)
\begin{itemize}
  \item Variantní text pro povolení s deklarací: ``Použití generativní AI je v tomto předmětu povoleno pouze se souhlasem vyučujícího a za podmínky, že při odevzdání student explicitně uvede: (1) které nástroje použil, (2) konkrétní prompty a verze, (3) vlastní úpravy a reflexi, a (4) přiloží pracovní verze/drafty.'' (general background knowledge)
  \item Variantní text pro zakázání: ``Použití externích generativních nástrojů k vytvoření obsahu zadání je zakázáno. Podezření na porušení integritní politiky bude řešeno podle akademických pravidel VUT.''
  \item Pole deklarace na odevzdání: krátké checkboxy a textové pole se záznamem: nástroj, prompt, popis přidané hodnoty (např. 150–300 znaků).
\end{itemize}

Rubrika hodnocení obsahující AI‑deklaraci — návrh polí (general background knowledge)
\begin{itemize}
  \item Transparentnost použití AI (0–5 bodů): student popíše rozsah použití AI a doloží prompty/konkrétní výstupy.
  \item Originalita a vlastní přínos (0–10 bodů): posuzuje se míra vlastní práce, úpravy a kritické evaluace AI‑výstupů.
  \item Doklad procesu (0–5 bodů): přiložené drafty, commity v repozitáři, screencasty či živá obhajoba.
\end{itemize}

Postupy při podezření z nevhodného použití (general background knowledge)
\begin{enumerate}
  \item Nejprve vyžádejte procesní důkazy (verze, drafty, artefakty) a krátkou obhajobu studenta.
  \item Proveďte krátké interview/ústní ověření, zaměřené na konkrétní kroky řešení.
  \item Pokud podezření přetrvává, zahajte interní proces podle akademických předpisů fakulty (dokumentovat průběh a rozhodnutí).
\end{enumerate}

Poznámka o detektorech a alternativách (general background knowledge)
\begin{itemize}
  \item Automatické „detektory AI“ mohou mít omezení a falešné pozitivy; místo spoléhání se pouze na ně kombinujte procesní ověřování, redesign úloh a požadavek transparentnosti v odevzdání.
\end{itemize}

Krátké doporučení pro implementaci na úrovni předmětu
\begin{enumerate}
  \item Na začátku semestru: aktualizujte sylabus podle jedné z výše uvedených kategorií a komunikujte očekávání studentům.
  \item Upravte rubriky tak, aby hodnotily proces i produkt a zahrnovaly pole pro AI‑deklaraci.
  \item Nasazujte postupy ověřování (drafty, ústní kontrola) u klíčových milníků.
  \item Zapojte metodické oddělení a PR/PRIV týmy při pilotním zavádění nových pravidel (zejména při zpracování PII nebo změnách v systému hodnocení).
\end{enumerate}

Závěrem. Cílem není zákaz technologií, ale udržení férovosti hodnocení a rozvoj dovedností studentů v kritickém používání AI. Implementace politiky by měla být praktická, srozumitelná a komunikovaná již v sylabu; kombinace procesních požadavků, deklarací a místních ověření je nejúčinnějším přístupem pro většinu předmětů (general background knowledge).